% Part: lambda-calculus
% Chapter: introduction
% Section: primitive-recursion

\documentclass[../../../include/open-logic-section]{subfiles}

\begin{document}

\olfileid{lam}{int}{pr}
\olsection{الدوال \usetoken{S}{lambda definable} مغلقة تحت العودية البدائية}

حين نصل إلى العودية البدائية، نحتاج أخيرًا إلى بعض العمل. وسنمضي على
مراحل. وكما سبق، نفترض أن لدينا حدين~$G'$ و$H'$، وأن كلًا منهما
!!{lambda define} دالته، $g$ و$h$ على الترتيب، ونريد حدًا~$F$
!!{lambda define}s الدالة~$f$ المعرّفة بـ
\begin{align*}
f(0, \vec z) & = g(\vec z) \\
f(x+1, \vec z) & = h(x, f(x,\vec z), \vec z).
\end{align*}
وبوجه عام، إذا أعطينا حدي لامبدا~$G'$ و$H'$، فيكفي إيجاد حد~$F$ بحيث
\begin{align*}
F(\num{0}, \vec z) & \equiv G(\vec z) \\
F(\overline{n+1}, \vec z) & \equiv H(\num{n}, F(\num{n}, \vec z), \vec z)
\end{align*}
لكل عدد طبيعي~$n$؛ وكون كل من~$G'$ و$H'$ !!{lambda define} دالته،
$g$ و$h$ على الترتيب، يعني أنه متى أحللنا الأعداد~$\num{\vec m}$ محل~$\vec z$،
اختزل~$F(\num{n+1}, \num{\vec m})$ إلى الجواب الصحيح في صورة سوية.

لكن يكفي لهذا إيجاد حد~$F$ يحقق
\begin{align*}
F(\num 0) & \equiv G \\
F(\overline {n+1}) & \equiv H(\num{n},F(\num{n}))
\intertext{لكل عدد طبيعي $n$، حيث}
G & = \lambd[\vec z][G'(\vec z)] \text{ و}\\
H(u,v) & = \lambd[\vec z][H'(u,v(\vec z),\vec z)].
\end{align*}
وبعبارة أخرى، نستطيع بحيل لامبدا تجنب القلق بشأن المعلمات الإضافية
$\vec z$، إذ يمتصها ترميز لامبدا.

وقبل تعريف الحد~$F$، نحتاج إلى آلية لمعالجة الأزواج المرتبة. وتوفرها
اللمّة الآتية.

\begin{lem}
يوجد حد لامبدا~$D$ بحيث يكون، لكل زوج من حدي لامبدا~$M$ و$N$،
$D(M,N)(\num{0}) \red M$ و$D(M,N)(\num{1}) \red N$.
\end{lem}

\begin{proof}
عرّف أولًا حد لامبدا~$K$ بـ
\[
K(y) = \lambd[x][y].
\]
وبعبارة أخرى،~$K$ هو الحد~$\lambd[y][\lambd[x][y]]$. ومن وجهة أخرى،
لكل~$M$، تكون~$K(M)$ دالة ثابتة تعيد~$M$ على أي مدخل.

وعرّف الآن~$D(x,y,z)$ بـ$D(x,y,z) = z (K(y))x$. عندئذ لدينا
\begin{align*}
D(M,N,\num 0) & \red \num 0 (K(N)) M \red M \text{ و}\\
D(M,N,\num 1) & \red \num 1 (K(N)) M \red K(N) M \red N,
\end{align*}
كما هو مطلوب.
\end{proof}

الفكرة هي أن~$D(M,N)$ يمثل الزوج~$\tuple{M, N}$، وإذا افترضنا أن~$P$
يمثل زوجًا كهذا، فإن~$P(\num 0)$ و$P(\num 1)$ يمثلان الإسقاطين الأيسر
والأيمن، أي~$(P)_0$ و$(P)_1$. وسنستعمل الترميزين الأخيرين.

\begin{lem}
الدوال~!!{lambda definable} مغلقة تحت العودية البدائية.
\end{lem}

\begin{proof}
نحتاج إلى إثبات أنه إذا أعطينا أي حدين~$G$ و$H$، أمكننا إيجاد حد~$F$
بحيث
\begin{align*}
F(\num{0}) & \equiv G \\
F(\num{n+1}) & \equiv H(\num{n}, F(\num{n}))
\end{align*}
لكل عدد طبيعي~$n$. والفكرة، على وجه التقريب، هي حساب متتاليات من
\emph{الأزواج}
\[
\tuple{\num 0, F(\num 0)}, \tuple{\num 1, F(\num 1)}, \dots,
\]
باستعمال الأعداد بوصفها مكررات. ولاحظ أن الزوج الأول ليس إلا
$\tuple{\num 0, G}$. وإذا أعطينا زوجًا~$\tuple{\num{n}, F(\num{n})}$،
فمن المفترض أن يكون الزوج التالي~$\tuple{\num{n+1}, F(\num{n+1})}$
مكافئًا للزوج~$\tuple{\num{n+1}, H(\num{n}, F(\num{n}))}$. وسنصمم حد
لامبدا~$T$ ينجز هذا الانتقال ذي الخطوة الواحدة.

والتفاصيل كما يأتي. عرّف~$T(u)$ بـ
\[
T(u) = \tuple{\fn{Succ}((u)_0), H((u)_0,(u)_1)}.
\]
ويسهل الآن التحقق من أنه، لكل عدد~$n$،
\[
T(\tuple{\num{n}, M}) \red \tuple{\num{n+1}, H(\num{n}, M)}.
\]
وكما اقترحنا، عند إعطاء~$G$ و$H$، عرّف~$F(u)$ بـ
\[
F(u) = (u(T,\tuple{\num 0, G}))_1.
\]
وبعبارة أخرى، عند المدخل~$\num{n}$، يكرر~$F$ الدالة~$T$ عدد~$n$ من
المرات على~$\tuple{\num 0, G}$، ثم يعيد المركبة الثانية. وفي البداية
لدينا
\begin{enumerate}
\item $\num{0} (T, \tuple{\num 0, G}) \equiv \tuple{\num{0}, G}$
\item $F(\num{0}) \equiv G$
\end{enumerate}
وبالاستقراء على~$n$، نستطيع إثبات أن الآتي يصدق لكل عدد طبيعي:
\begin{enumerate}
\item $\num{n+1}(T, \tuple{\num{0}, G}) \equiv \tuple{\num{n+1}, F(\num{n+1})}$
\item $F(\num{n+1}) \equiv H(\num{n}, F(\num{n}))$
\end{enumerate}
فللبند الثاني لدينا
\begin{align*}
F(\num{n+1}) & \red (\num{n+1}(T, \tuple{\num 0, G}))_1 \\
& \equiv (T(\num{n} (T, \tuple{\num 0, G})))_1 \\
& \equiv (T(\tuple{\num{n}, F(\num{n})}))_1 \\
& \equiv (\tuple{\num{n+1}, H(\num{n}, F(\num{n}))})_1 \\
& \equiv H(\num{n}, F(\num{n})).
\end{align*}
وقد استعملنا هنا فرض الاستقراء في السطر قبل الأخير. وللبند الأول لدينا
\begin{align*}
\num{n+1} (T, \tuple{\num 0, G}) &
\equiv T(\num{n} (T, \tuple{\num 0, G})) \\
& \equiv T( \tuple{\num{n}, F(\num{n})}) \\
& \equiv \tuple{\num{n+1}, H(\num{n}, F(\num{n}))} \\
& \equiv \tuple{\num{n+1}, F(\num{n+1})}.
\end{align*}
وقد استعملنا هنا البند الثاني في السطر الأخير. وهكذا أثبتنا أن
$F(\num 0) \equiv G$ وأنه، لكل~$n$، يتحقق~$F(\num {n+1}) \equiv H(\num n,
F(\num{n}))$، وهو بالضبط ما نحتاج إليه.
\end{proof}

\end{document}
